% Part: sets-functions-relations
% Chapter: functions
% Section: partial-functions

\documentclass[../../../include/open-logic-section]{subfiles}

\begin{document}

\olfileid{sfr}{fun}{par}

\olsection{偏函数}

\begin{explain}
有时放宽函数的定义是有用的，即不要求函数的输出对所有可能的输入都有定义。
这样的映射称为\emph{偏函数}。
\end{explain}

\begin{defn}
一个\emph{偏函数} $f \colon A \pto B$ 是一个映射，它将 $B$ 中至多一个元素指派给 $A$ 的每个元素。
如果 $f$ 将 $B$ 的一个元素指派给 $x \in A$，我们就说 $f(x)$ 是\emph{有定义的}，否则就是\emph{未定义的}。
如果 $f(x)$ 有定义，我们记作 $f(x) \fdefined$，否则记作 $f(x) \fundefined$。
偏函数~$f$ 的\emph{定义域}是 $A$ 中使得 $f$ 有定义的子集，即$\dom{f} = \Setabs{x \in A}{f(x) \fdefined}$。
\end{defn}

\begin{ex}
每个函数 $f\colon A \to B$ 也是一个偏函数。那些在 $A$ 上处处有定义的偏函数——也就是我们之前简称为函数的那种——也称为\emph{全函数}。
\end{ex}

\begin{ex}
由 $f(x) = 1/x$ 给出的偏函数 $f \colon \Real \pto \Real$ 在 $x = 0$ 处未定义，在其他各处都有定义。
\end{ex}

\begin{prob}
给定 $f\colon A \pto B$，如下定义偏函数 $g\colon B \pto A$：对任意 $y \in B$，如果存在唯一的 $x \in A$ 使得 $f(x) = y$，则令 $g(y) = x$；否则令 $g(y) \fundefined$。试证：若 $f$ 是单射，则对所有 $x \in \dom{f}$ 有 $g(f(x)) = x$，且对所有 $y \in \ran{f}$ 有 $f(g(y)) = y$。
\end{prob}

\begin{defn}[偏函数的图像]
设 $f\colon A \pto B$ 是一个偏函数。$f$ 的\emph{图像}是关系 $R_f \subseteq A \times B$，定义为
\[
R_f = \Setabs{\tuple{x,y}}{f(x) = y}.
\]
\end{defn}

\begin{prop}
设 $R \subseteq A \times B$ 满足如下性质：只要 $Rxy$ 且 $Rxy'$，就有 $y = y'$。那么 $R$ 就是如下定义的偏函数 $f\colon A \pto B$ 的图像：若存在 $y$ 使得 $Rxy$，则令 $f(x) = y$，否则令 $f(x) \fundefined$。如果 $R$ 还是\emph{序列的}，即对每个 $x \in A$ 都存在 $y \in B$ 使得 $Rxy$，那么 $f$ 就是全函数。
\end{prop}

\begin{proof}
假设存在 $y$ 使得 $Rxy$。如果还存在另一个 $y' \neq y$ 使得 $Rxy'$，则违背了关于 $R$ 的条件。因此，如果存在 $y$ 使得 $Rxy$，那么这个 $y$ 是唯一的，所以 $f$ 是良定义的。显然，$R_f = R$，并且若 $R$ 是序列的，则 $f$ 是全函数。
\end{proof}

\end{document}